\section{Apply Reverse AD to a \emph{map - reduce min} Composition}
My code validates and is provided below:
\begin{lstlisting}[language={futhark}]
def task_mapomin [n] (as: [n]f32) (y_bar: f32) : (f32, [n]f32) =
  let bs = map (\a -> let b = a * a - 0.5 * a in b) as
  let (imin, y) =
    zip (iota n) bs
    |> reduce_comm (\(i1, b1) (i2, b2) ->
                      if b1 == b2
                      then if i1 < i2
                           then (i1, b1)
                           else (i2, b2)
                      else if b1 < b2
                      then (i1, b1)
                      else (i2, b2))
                   (n, f32.highest)
  let bs_bar = tabulate n (\i -> if i == imin then y_bar else 0.0f32)
  let as_bar =
    map2 (\a b_bar ->
            let a_bar = (2 * a - 0.5) * b_bar
            in a_bar)
         as
         bs_bar
  in (y, as_bar)
\end{lstlisting}

The lifted operator is what is the function passed to the reduce\_comm
function. This is both associative and cumulative, as I do tie breaks on the
indices. Essentially this works like the regular minimum function, but breaking
ties by choosing the smaller of the two indices. This means that no matter what
order you give the inputs, if the inputs are equivalent, then the result is the
same.

First, then primal is called, where we have substituted the reduce\_comm call
with the one including the lifted operator mentioned before. Next, we have the
tabulate call, which is what corresponds to the adjoint of the reduce\_comm.
This tabulate also populates the $\bar{as}$, which previously is not used and
as such is initialized to 0. The last map, then corresponds to the adjoint of
the map.

\begin{table}[H]
\centering
\caption{Benchmark runtimes and speedups for MapOmin.}
\label{tab:min}
\begin{tabular}{l cc cc}
\toprule
Method & \multicolumn{2}{c}{Runtime (\textmu s)} & \multicolumn{2}{c}{Overhead from primal} \\
\cmidrule(lr){2-3} \cmidrule(lr){4-5}
       & 10M elements & 100M elements & 10M elements & 100M elements \\
\midrule
pomin\_primal & 67  & 303   & -    & - \\
vjp2          & 157 & 910   & 2.34 & 3.00 \\
manual        & 147 & 1029  & 2.19 & 3.40 \\
\bottomrule
\end{tabular}
\end{table}

The runtimes can be seen in Table~\ref{tab:min}. Here we see that in general
the run times of my implementation is very close to that of the built-in vjp2,
which assuming it also does the same optimisation is to be expected.
